% Pursuing Stacks (cotes 134-2 to 134-8): archive cote 134-6, site page 2 (PDF page 3), typescript page 321.
% typescript. Draft compared: ps5.tex:1412-1477.
% AI-transcribed by Claude Opus 5.5 (claude-opus-5-5, adaptive thinking, effort high), Message Batches API, batch ps3 (msgbatch_01XdFmuVVBAcgc7EFxdtCqsj), 2026-11-12; not yet checked by eye.
\scanpage{134-6/2}
\mspage{321}

\margin{\{}topos (with suitable assumptions on A, M ...).
d) Possibility of expressin\over{g}{g} finite type of a homotopy type in terms
of $A^\wedge_{fp}$ , for suitable test categories\hand{A}.

\margin{$\uparrow$}e) Possibility of defining boundary ope\xout{\illegible{1}}ations within a test cate-%
gory - and/or getting Dold-Puppe type relations.

\hand{* * *}

\margin{$\downarrow$}Before resuming more technical work with the matters left over for
part IV, I would like still to write down some afterthoughts, concer-%
ning the question of boundary operations in a test category (question
8) in our review). It occured to me that perhaps it is'nt such a good
idea, to try at all costs to subordinate this question to a question of
cellular decompositions of spheres, however natural this idea may be
in view of the example of the standar\over{t}{d} types of complexes and multicom-%
plexes. In this connection, I remember that among my first thoughts
when starting unwittingly on the modelizing story, was that a ``test
category'' A (namely one such that $A^\wedge$ should be ``modelizing'' ) should
more or less correspond to such \deleted{subdivisions}\inserted{decompositions}. Soonafter it came as a
big surprise that so little was needed in fact for A to merit the name
of a test category - and that the \xout{\illegible{1}} relevant conditions had nothing
to do with cellular decompositions of this or that. T\over{h}{h}e same may well
turn out, when looking for a generalization of the standard simplicial,
cubical or hemispherical \inserted{chain} complexes, giving rise to the homology and
cohomology invariants of a given ``complex''. T\over{h}{h}e kind of set up I pro-%
posed in yesterday's notes, for a formalism of boundary operations
in a test category A, now looks to me in some respects somewhat \xout{"etriqu}
``étriqué'' , and I'll try another start in a different spirit.

In order not to get involved in irrelevant technicalities, I as-%
sume that the basic localizer $\underline{W}$ is $\underline{W}_\infty$ = usu\xout{\illegible{1}}al weak equivalence. It
seems that \xout{at} one basic fact \xout{\illegible{2}} for writing down relationship
between homotopy types and ``homology types'' , is the existence of a
canonical ``abelianization functor''
\[
\text{(1)}\qquad (\mathrm{Hot}) \longrightarrow \mathrm{D.}((\mathrm{Ab}))\hand{$\Bigl(\overset{\mathrm{def}}{=}\ \mathrm{Hotab}\Bigr)$,}
\]
where (Ab) is the abelian category of abelian groups, and D. designa\xout{\illegible{1}}
tes the ``derived category'' of \xout{an abelian category} the category
Ch.((Ab)) of chain complexes of abelian groups, namely its localization
with respect to\hand{``}weak equivalences\over{,}{''} \hand{i.e. quasi-isomorphisms:}

\margin{2}
